\documentclass[a4paper]{article}

\usepackage[margin=1cm]{geometry}
\usepackage{CJKutf8}
\usepackage[utf8]{inputenc}
\usepackage{tikz}
\usetikzlibrary{arrows.meta, positioning, shapes.geometric, calc}
\pagestyle{empty}

\begin{document}
\begin{CJK*}{UTF8}{gbsn}

\begin{center}
{\Large\bfseries GNS-J1900 内置 U 盘镜像 更新与打包 流程图}
\end{center}
\vspace{3mm}

\resizebox{\textwidth}{!}{%
\begin{tikzpicture}[
  >=Stealth,
  font=\scriptsize,
  file/.style   = {rectangle, draw=blue!55!black, thick, rounded corners=2pt,
                    fill=blue!8, text width=4.2cm, align=center,
                    minimum height=1.3cm, inner sep=3pt},
  proc/.style   = {rectangle, draw=orange!70!black, thick, rounded corners=2pt,
                    fill=orange!12, text width=4.2cm, align=center,
                    minimum height=1.3cm, inner sep=3pt},
  cfg/.style    = {rectangle, draw=teal!60!black, thick, rounded corners=2pt,
                    fill=teal!8, text width=4.2cm, align=center,
                    minimum height=1.1cm, inner sep=3pt},
  deliver/.style= {rectangle, draw=green!45!black, very thick, rounded corners=2pt,
                    fill=green!12, text width=4.2cm, align=center,
                    minimum height=1.4cm, inner sep=3pt, font=\scriptsize\bfseries},
  merge/.style  = {circle, draw=black, thick, fill=yellow!35, minimum size=7mm,
                    inner sep=0pt, font=\bfseries},
  lbl/.style    = {font=\tiny, midway, above, text width=2.6cm, align=center},
]

% ---------- 列坐标: col1=0  col2=6.2  col3=12.4 ----------

% ===== 子流程：配置自动生成（位于最上方） =====
\node[file] (dir) at (6.2, 9.4) {扫描目录\\ \texttt{DIR}};
\node[proc] (pygen) at (6.2, 7.1) {\texttt{gns\_iso\_to\_builtin\_}\\\texttt{u\_disk\_image.py}\\\texttt{-g DIR}\ 扫描生成};
\node[cfg]  (toml) at (6.2, 4.8) {自动生成\\\texttt{gns\_iso\_to\_builtin\_}\\\texttt{u\_disk\_image.toml}};

\draw[->, thick] (dir) -- (pygen);
\draw[->, thick] (pygen) -- (toml);

% ===== 输入：旧镜像 / 新 ISO（第一行左端，纵向堆叠） =====
\node[file] (old) at (0, 7.1) {旧的\\\texttt{gns-j1900-internal-}\\\texttt{u-disk-YYYYYYYY.img.gz}};
\node[file] (newiso) at (0, 4.8) {新的\\\texttt{GNS-X-X32-X-}\\\texttt{Debian12(64)-Oto35V20-}\\\texttt{STD001-XXXXXXXX.iso}};

% ===== 第一行（从左到右）：合并 -> 主处理 -> 输出镜像 =====
\node[merge] (m1) at (0, 2.0) {+};
\draw[->, thick] (old.west) -- ($(old.west)+(-1.2,0)$) |- ($(m1.west)+(-1.2,0)$) -- (m1.west);
\draw[->, thick] (newiso) -- (m1);

\node[proc] (p1) at (6.2, 2.0) {\texttt{gns\_iso\_to\_builtin\_}\\\texttt{u\_disk\_image.py}\\(主处理)};
\draw[->, very thick] (m1) -- (p1);
\draw[->, thick] (toml) -- (p1) node[lbl] {提供配置};

\node[deliver] (newimg) at (12.4, 2.0) {(a) 交付\\\texttt{gns-j1900-internal-}\\\texttt{u-disk-XXXXXXXX.img.gz}};
\draw[->, very thick] (p1) -- (newimg);

% ===== 蛇形转向：第一行末端(col3) 向下折到第二行(col3)，再向左 =====
\node[merge] (m3) at (12.4, -2.6) {+};
\draw[->, very thick] (newimg) -- (m3);

\node[file] (empty) at (12.4, -5.4) {空白 ISO 模板\\\texttt{GNS-J1900-Internal-}\\\texttt{U-Disk-Empty.iso}};
\draw[->, very thick] (empty) -- (m3);

\node[deliver] (finaliso) at (6.2, -2.6) {最终交付\\\texttt{GNS-J1900-Internal-}\\\texttt{U-Disk-XXXXXX.iso}};
\draw[->, very thick] (m3) -- (finaliso);

\end{tikzpicture}%
}

\vspace{4mm}
\begin{flushleft}
\footnotesize
\textbf{说明：}\quad \par
(a) 更新内置 U 盘镜像：旧镜像 + 新 ISO 经 \texttt{gns\_iso\_to\_builtin\_u\_disk\_image.py} 处理（配置文件 \texttt{.toml} 由同一脚本以 \texttt{-g DIR} 扫描指定路径自动生成）得到新镜像；   \par
(b) 打包为 ISO：新镜像与空白 ISO 模板合并，最终生成 \texttt{GNS-J1900-Internal-U-Disk-XXXXXX.iso} 交付。
\end{flushleft}

\end{CJK*}
\end{document}
